\section{Introduction}
\label{sec:introduction}

Courses, editorial teams and small organisations decide every day which documents deserve
further attention. Manual review is accurate but slow, because the reviewer must read and
record the evidence; a purely automatic classifier is fast but hides the reason for its
decision inside the model. When the decision affects authors, partners or the allocation of
scarce attention, the reviewer needs more than a label: the evidence behind the label must
remain inspectable.

This report asks: \emph{in human-in-the-loop document triage, can we shorten review time
while keeping the human decision and its supporting evidence inspectable?} We answer with a
reproducible prototype rather than a new classification model.

The contribution is threefold:
\begin{itemize}
  \item a three-stage review workflow: documents enter, the model proposes a label and
        evidence, and a human accepts, edits or rejects the proposal;
  \item a reproducible prototype with a pinned environment, a synthetic corpus, an audit
        trail and documented failure cases;
  \item a small internal evaluation reporting review time, override rate and evidence
        completeness, together with the limits of that evaluation.
\end{itemize}

Section~\ref{sec:related} reviews related work, Section~\ref{sec:method} describes the
method, Section~\ref{sec:implementation} covers implementation and reproducibility,
Section~\ref{sec:evaluation} reports the evaluation, Section~\ref{sec:discussion} discusses
ethics, accountability and limitations, and Section~\ref{sec:conclusion} concludes.
